
\begin{abstract} \NS
La remediación del agua se ha convertido en un tema de creciente interés en todo el mundo debido a la necesidad de recuperar y reutilizar las aguas residuales en muchos sectores, como la agricultura y el reciclaje industrial. La tecnología del plasma no térmico es una técnica prometedora para la remediación del agua y se ha convertido en un campo de intensa investigación. En el plasma no térmico se generan especies altamente reactivas como $\tx{O}$, $\tx{OH}$, $\tx{O}_3$, así como también radiación UV. Esta tecnología combina la contribución de especies activas y condiciones físicas que han demostrado una gran eficiencia en la degradación de muchos compuestos orgánicos, así como en la destrucción e inactivación de virus y bacterias. \\\\
En este trabajo se presenta el montaje de un reactor de plasma trielectródico para el tratamiento de aguas contaminadas. En el reactor, se genera una descarga de barrera dieléctrica (DBD) en el aire ambiente mediante la aplicación de alta tensión alterna entre dos electrodos planos fijados a lados opuestos de un material dieléctrico. A una distancia de unos pocos centímetros, se fija a un canal por el que circula el agua a tratar un tercer electrodo conectado a tierra. Al elevar el potencial del sistema de electrodos DBD con un alto voltaje continuo, los microcanales de plasma no térmico generados en la DBD se propagan hacia el tercer electrodo y alcanzan la superficie del agua. \\\\
Este diseño es escalable ya que permite la conexión en paralelo de varios sistemas de electrodos DBD para la generación de múltiples descargas sucesivas a lo largo del canal por el cual fluye el agua a tratar. El funcionamiento del reactor, en términos de eficiencia de eliminación y rendimiento energético, se evaluó mediante el tratamiento de una solución acuosa de azul de metileno. Se presentan también resultados preliminares combinando el tratamiento con plasma con recubrimientos de $\tx{TiO}_2$ ubicados sobre el canal de circulación del agua en la zona de la descarga.
\end{abstract}